\title[Causality]{Causality: Lecture 10}
\author[Joris Mooij]{\\[5mm]\large Joris Mooij \& Patrick Forr\'e\\
\url{j.m.mooij@uva.nl}, \url{p.d.forre@uva.nl}}
\institute[UvA]{University of Amsterdam}
\date[2022-04-11]{\normalsize April 11th, 2022}
\maketitle


\part{Equivalence Relations}
\frame{\partpage}

\begin{frame}{Equivalence relations}
Equivalence relations capture the notion that mathematical objects can be ``equivalent'' from some point of view.
\begin{definition}
  Let $Z$ be a set and $R \subseteq Z^2$ be a relation on $Z$ (i.e., a subset of ordered pairs of $Z$). $R$ is called an \emph{equivalence relation} if 
  \begin{enumerate}
    \item $R$ is reflexive: $(a,a) \in R$ for all $a \in Z$;
    \item $R$ is symmetric: $(a,b) \in R \iff (b,a) \in R$ for all $a,b \in Z$;
    \item $R$ is transitive: if $(a,b) \in R$ and $(b,c) \in R$ then $(a,c) \in R$ for all $a,b,c\in Z$.
  \end{enumerate}
\end{definition}
\end{frame}

\begin{frame}{Equivalence of SCMs}
\begin{definition}
  Let $M = \lt \Sin, \Send, \Srnd, \CX, P, f \rt$ and
  $\tilde{M} = \ltuple \Sin, \Send, \Srnd, \CX, P, \tilde{f} \rtuple$ be two SCMs.
  We say that
  $M$ is \emph{equivalent} to $\tilde{M}$ and write $M \equiv \tilde{M}$ if for each $v \in \Send$,
  $$\forall x \in \CX: \quad x_v = f_v(\xJ,\xV,\xW) \iff x_v = \tilde{f}_v(\xJ,\xV,\xW).$$
%  In words, $M$ and $\tilde{M}$ are considered equivalent if they at most differ in terms of their
%  causal mechanisms, yet each of their structural equations has the same solutions.
\end{definition}
\begin{block}{Proposition}
\begin{enumerate}
  \item Equivalence of SCMs is preserved by hard interventions.
  \item Unique solvability and simplicity are invariants of SCM equivalence.
  \item Equivalent SCMs have the same solution functions, solutions, (potential) outcomes, and Markov kernels.
\end{enumerate}
\end{block}
\end{frame}

\begin{frame}{Example of equivalent SCMs}
\begin{example}
  The SCM with the single structural equation
  $$X = -X^3 + X + W$$
  where $X$ is endogenous, and $W$ is exogenous, is equivalent to the SCM obtained by replacing the structural equation with
  $$X = \sqrt[3]{W}.$$
  Note that $X$ no longer appears on the r.h.s..
\end{example}
\end{frame}

\begin{frame}{Other equivalences: definition}
\begin{definition}
  Let $M = \lt \Sin, \Send, \Srnd, \CX, P, f \rt$ and
  $\tilde{M} = \ltuple \tilde{\Sin}, \tilde{\Send}, \tilde{\Srnd}, \tilde{\CX}, \tilde{P}, \tilde{f} \rtuple$ be two simple SCMs
  and $O \subseteq (\Send \cap \tilde{\Send}) \cup (\Srnd \cap \tilde{\Srnd})$ a subset.
  We say that:
  \begin{enumerate}
    \item $M$ and $\tilde{M}$ are \emph{observationally equivalent w.r.t.\ $O$} if $\CX_O = \tilde{\CX}_O$, $\CX_{J \cap \tilde{J}} = \tilde{\CX}_{J \cap \tilde{J}}$ and their marginal Markov kernels coincide:
      $$P_{M}(X_O \mid \doit(X_J)) = P_{\tilde{M}}(X_O \mid \doit(X_{\tilde J})).$$
    This has to be interpreted as both Markov kernels being a version of a Markov kernel $\CX_{J \cap \tilde{J}} \dshto \CX_O$,
    i.e., $P_{M}(X_O \mid \doit(X_J))$ must be essentially constant in $x_{J \setminus \tilde{J}}$ and
    $P_{\tilde{M}}(X_O \mid \doit(X_{\tilde{J}}))$ must be essentially constant in $x_{\tilde{J} \setminus J}$;
%    \item $M$ and $\tilde{M}$ are \emph{observationally equivalent w.r.t.\ $O$} if $\CX_O = \tilde{\CX}_O$, $J = \tilde{J}$, $\CX_J = \tilde{\CX}_J$ and $P_{M}(X_O \mid \doit(\XJ)) = P_{\tilde{M}}(X_O \mid \doit(\XJ))$;
    \item $M$ and $\tilde{M}$ are \emph{interventionally equivalent w.r.t.\ $O$} if for every subset $T \subseteq O$ the intervened SCMs $M_{\doit(T)}$ and $\tilde{M}_{\doit(T)}$ are observationally equivalent w.r.t.\ $O \setminus T$.
%    \item $M$ and $\tilde{M}$ are said to be \emph{counterfactually equivalent w.r.t.\ $O$} if the twin SCMs $M^\twin$ and $\tilde{M}^\twin$ are interventionally equivalent w.r.t.\ $O \cup O'$, where $O'$ is the copy of $O$ in $\Send' \cap \tilde{\Send}'$.
  \end{enumerate}
\end{definition}
\end{frame}

\begin{frame}{Other equivalences: relationships}
\begin{block}{Proposition}
  For simple SCMs $M, \tilde{M}$ and a subset $O \subseteq (\Send \cap \tilde{\Send}) \cup (\Srnd \cap \tilde{\Srnd})$:
  \begin{enumerate}
%    \item If $M \equiv \tilde{M}$ then $M$ and $\tilde{M}$ are counterfactually equivalent w.r.t.\ $O$.
    \item If $M$ and $\tilde{M}$ are equivalent, then 
      $M$ and $\tilde{M}$ are interventionally equivalent w.r.t.\ $O$.
    \item If $M$ and $\tilde{M}$ are interventionally equivalent w.r.t.\ $O$ then 
      $M$ and $\tilde{M}$ are observationally equivalent w.r.t.\ $O$.
  \end{enumerate}
\end{block}
The reverse implications do not hold in general. This expresses that causal modeling is more refined than probabilistic modeling.
\end{frame}

\begin{frame}{Counterexample}
\begin{example}
  The following SCMs are both simple:
  \begin{align*}
    M &: \qquad N \sim \C{N}(\mu,\sigma^2) \qquad C = \alpha + \beta N \\
    \tilde{M}&: \qquad C \sim \C{N}(\tilde{\mu},\tilde{\sigma}^2) \qquad N = \tilde{\alpha} + \tilde{\beta} C
  \end{align*}
Their observational distributions are respectively 
  $$P_M\left(\begin{pmatrix}N \\C\end{pmatrix}\right) = \C{N}\left(\begin{pmatrix} \mu \\ \alpha + \beta \mu \end{pmatrix}, \begin{pmatrix} \sigma^2 & \beta\sigma^2 \\ \beta\sigma^2 & \beta^2\sigma^2\end{pmatrix}\right)$$
and
  $$P_{\tilde{M}}\left(\begin{pmatrix}N\\C\end{pmatrix}\right) = \C{N}\left(\begin{pmatrix} \tilde{\alpha} + \tilde{\beta}\tilde{\mu} \\ \tilde{\mu} \end{pmatrix}, \begin{pmatrix} \tilde{\beta}^2\tilde{\sigma}^2 & \tilde{\beta}\tilde{\sigma}^2 \\ \tilde{\beta}\tilde{\sigma}^2 & \tilde{\sigma}^2\end{pmatrix}\right)$$
For certain parameter choices, they are observationally equivalent.
They are NOT interventionally equivalent except for very special parameter choices.
\end{example}
\end{frame}

\part{Marginalizations}
\frame{\partpage}

\begin{frame}{Marginalizations: Motivation}
  Often we like to obtain a simpler description of a system by hiding irrelevant details of a subsystem.

  \begin{block}{Examples}
    \begin{enumerate}
      \item Moving code into a subroutine when programming
      \item Marginalizing a probability distribution over some variables
      \item Taking an orthogonal projection in linear algebra
      \item Taking a latent projection of a CDMG
    \end{enumerate}
  \end{block}

  For SCMs, we define a marginalization over a subset of endogenous variables that removes those
  variables from the model but preserves the causal semantics of the remaining variables.
\end{frame}

\begin{frame}{Marginalizations: Basic insight}
  Write $K := V \setminus L$ and let $\tilde{g}_L$ denote the unique solution function of $M_{\doit(K)}$.
    \begin{equation*}\begin{split}
%  & \begin{cases}
%    x_L = g_L(x_J,x_W) \\
%    x_M = g_M(x_J,x_W) \\
%  \end{cases} \\
%  \iff 
  & \begin{cases}
    x_L = f_L(x_J,x_K,x_L,x_W) \\
    x_K = f_K(x_J,x_K,x_L,x_W) \\
  \end{cases} \\
  \iff & \begin{cases}
    x_L = \tilde{g}_L(x_J,x_K,x_W) \\
    x_K = f_K(x_J,x_K,x_L,x_W) \\
  \end{cases}  \\
  \iff & \begin{cases}
    x_L = \tilde{g}_L(x_J,x_K,x_W) \\
    x_K = f_K(x_J,x_K,\tilde{g}_L(x_J,x_K,x_W),x_W) =: \tilde{f}(\xJ,x_{K},\xW) \\
  \end{cases} \\
%  \iff & \begin{cases}
%    x_L = \tilde{g}_L(x_J,x_K,x_W) \\
%    x_K = \tilde{f}(x_J,x_K,x_W) \\
%  \end{cases}
  \end{split}\end{equation*}
  We can now forget $x_L$, and only retain:
  $$x_K = \tilde{f}(x_J,x_K,x_W)$$
\end{frame}

\begin{frame}{Marginalizations: Definition}
\begin{Def}
  Let $L \subseteq \Send$ be such that $M_{\doit(V\setminus L)}$ is uniquely solvable.
  Write $K = V \setminus L$ and let $\tilde{g}_L : \CX_{\Sin \cup K} \times \CXW \to \CX_L$
  be the unique solution function for $M_{\doit(K)}$.
  Then we call $M_{\setminus L} = \ltuple \Sin, K, \Srnd, \CX_{\Sin} \times \CX_{K} \times \CX_{\Srnd}, P, \tilde{f} \rtuple$ with
  $$\tilde{f}(\xJ,x_{K},\xW) = f_{K}(\xJ,x_{K},\tilde{g}_L(\xJ,x_{K},\xW),\xW)$$
  the \emph{marginalization of $M$ over $L$}.
\end{Def}
The marginalization is obtained by performing the following steps:
  \begin{enumerate}
    \item Solve the submodel $M_{\doit(K)}$ on $L$
    \item Substitute the solution function into the causal mechanism for $K$
    \item Restrict to the $K$ component only
  \end{enumerate}
\end{frame}

\begin{frame}{Marginalization: Properties I}
In the next two slides, assume that
$L \subseteq \Send$ is such that $M_{\doit(V \setminus L)}$ is uniquely solvable
(i.e., $M_{\setminus L}$ is defined).
\begin{Lem}
  If $M$ is uniquely solvable then 
  \begin{itemize}
    \item its marginalization $M_{\setminus L}$ is uniquely solvable; 
    \item the unique solution function for $M_{\setminus L}$ is just $g_{V\setminus L}$, where 
  $g : \CX_J \times \CX_W \to \CX_V$ is the unique solution function for $M$.
  \end{itemize}
\end{Lem}

This gives:
$$P_{M_{\setminus L}}(X_{V\setminus L},X_W \mid \doit(X_J)) = P_{M}(X_{V\setminus L},X_W \mid \doit(X_J))$$
(The Markov kernel of the marginalized SCM is given by the marginal Markov kernel of the SCM).
\end{frame}
  
\begin{frame}{Marginalization: Properties II}
  \begin{Prp}
    For a hard intervention $\doit(T\dots)$ with target $T \subseteq \Sin \cup \Srnd \cup \Send$ (of any of the three variants) such that $L \cap T = \emptyset$:
    $$(M_{\doit(T\dots)})_{\setminus L} = (M_{\setminus L})_{\doit(T\dots)}.$$
\end{Prp}
Try proving this yourself!

\begin{Prp}
  Let $L_2 \subseteq \Send$ such that $L \cap L_2 = \emptyset$.
  If $(M_{\setminus L})_{\doit((V\setminus L)\setminus L_2)}$ is uniquely solvable, then $M_{\doit(V\setminus(L\cup L_2))}$ is uniquely solvable, and:
  $$(M_{\setminus L})_{\setminus L_2} = M_{\setminus (L\cup L_2)}.$$
\end{Prp}
\begin{proof}
  Straightforward, but cumbersome notation. See lecture notes.
\end{proof}
\end{frame}

\begin{frame}{Marginalization: Properties for Simple  SCMs}
Let $M = \lt \Sin, \Send, \Srnd, \CX, P, f \rt$ be a simple SCM.
\begin{Prp}
  For any $L \subseteq \Send$, the marginalization $M_{\setminus L}$ is also simple.
\end{Prp}
\begin{proof}
  For $T \subseteq J \cup W \cup V \setminus L$, marginalization commutes with hard intervention:
  $$(M_{\setminus L})_{\doit(T)} = (M_{\doit(T)})_{\setminus L}.$$
  Because $M$ is simple, also $M_{\doit(T)}$ is simple, and in particular it is
  uniquely solvable. 
  Hence its marginalization $(M_{\doit(T)})_{\setminus L}$ is uniquely solvable.
  This means that $(M_{\setminus L})_{\doit(T)}$ is uniquely solvable. Since this
  holds for any $T$, $M_{\setminus L}$ is simple.
\end{proof}
\begin{Cor}
  For $L_1, L_2 \subseteq \Send$ with $L_1 \cap L_2 = \emptyset$:
    $(M_{\setminus L_1})_{\setminus L_2} = (M_{\setminus L_2})_{\setminus L_1} = M_{\setminus (L_1 \cup L_2)}$.
\end{Cor}
\end{frame}

\begin{frame}{Marginalizations preserve causal semantics}
\begin{Thm}
  Let $M = \lt \Sin, \Send, \Srnd, \CX, P, f \rt$ be a simple SCM and $L \subseteq \Send$.
  Then $M$ and $M_{\setminus L}$ are observationally and
  interventionally %and counterfactually equivalent 
  w.r.t.\ $(V \cup W) \setminus L$.
\end{Thm}
\begin{proof}
  Write $K := V \setminus L$. We already noticed that
  $P_{M}(X_K,X_W \mid \doit(X_J)) = P_{M_{\setminus L}}(X_K,X_W \mid \doit(X_J))$.

  Let $T \subseteq K \cup W$. Then $(M_{\setminus L})_{\doit(T)} = (M_{\doit(T)})_{\setminus L}$.
  The observational equivalence of $M_{\doit(T)}$ and $(M_{\doit(T)})_{\setminus L}$ w.r.t.\ $(K \cup W) \setminus T$ hence implies the observational equivalence of $M_{\doit(T)}$ and $(M_{\setminus L})_{\doit(T)}$ w.r.t.\ $(K \cup W) \setminus T$.
  Hence, $M$ and $M_{\setminus L}$ are interventionally equivalent w.r.t.\ $K \cup W$.
\end{proof}
Note: this can be shown for general (non-simple) SCMs as well.
\end{frame}

\begin{frame}{Linear SCMs}
  \begin{Def}
We call an SCM $M=\lt \Sin, \Send, \Srnd, \CX, P, f \rt$
\emph{linear} if all variables are real-valued and
each component of the causal mechanism is an affine combination of endogenous variables, that is
$$
%  f_v(x) = \sum_{w\in V \cup J \cup W} B_{vw} x_w,
    f_v(x) = \sum_{u\in V} B_{vu}(x_J,x_W) x_u + c_v(x_J,x_W),
$$
    where $B(x_J,x_W)\in\RN^{V \times V}$ is a matrix-valued function $\CXJ\times\CXW \to \RN^{V\times V}$ and
$c(x_J,x_W) \in \RN^V$ is a vector-valued function $\CXJ\times\CXW \to \RN^V$.
  \end{Def}

  \begin{Prp}
Let $L \subseteq V$.
$(M)_{\doit(V\setminus L)}$ is uniquely solvable if and only if 
$I_{L}-B_{LL}(x_J,x_W)$ are invertible $\forall x_J \in \CX_J, x_W \in \CX_W$,
with unique solution function:
\small
  \begin{equation*}\begin{split}
    g &: \RN^{V\setminus L} \times \RN^{J\cup W} \to \RN^L \\
      &: (x_{V\setminus L},x_J,x_W) \mapsto (I_L - B_{LL}(x_J,x_W))^{-1} \big(B_{L,{V\setminus L}}(x_J,x_W) x_{V\setminus L} + c_{L}(x_J,x_W)\big).
  \end{split}\end{equation*}
  \end{Prp}
\end{frame}

\part{Graphs}
\frame{\partpage}

\begin{frame}{Finally, graphs!}
  We introduce the graph of an SCM with nodes corresponding to the variables and directed edges representing the functional relationships, i.e., the structure of the structural equations.
  
  \begin{Def}
  For $i \in \Sin \cup \Send \cup \Srnd$ and $j \in \Send$, we say that
  \emph{$i$ is a parent of $j$ according to $M$} if there does not exist a measurable function
  $\tilde f_j : \CX_{(\Sin\cup\Send\cup\Srnd) \setminus \{i\}} \to \CX_j$
  such that for all $x \in \CX$,
  $$x_j = f_j(x) \iff x_j = \tilde f_j(x_{\setminus i}),$$
  where $x_{\setminus i}$ is shorthand for $x_{(\Sin\cup\Send\cup\Srnd)\setminus\{i\}}$.
\end{Def}
  \begin{itemize}
    \item In words: $i$ is parent of $j$ if the causal mechanism for $j$ is not equivalent to one that is constant w.r.t.\ its input component $i$.
    \item By definition, exogenous (input and random) variables have no parents.
    \item Note that the parent-relationship is preserved under equivalence.
  \end{itemize}
\end{frame}

\begin{frame}{Graphs: Definition}
  \begin{Def}
  The CDG $\lt \Sin, \Send \cup \Srnd, E \rt$ with
  input nodes $\Sin$, output nodes $\Send\cup\Srnd$, and
  directed edges
  $$E = \{i \to j : i \in \Sin\cup\Srnd\cup\Send, j\in\Send: \text{$i$ is parent of $j$ according to $M$}\}$$
  is called the \emph{graph of the SCM $M$} and will be denoted as $G(M)$.
\end{Def}
  \begin{itemize}
    \item Equivalent SCMs have the same graphs.
    \item Observationally equivalent SCMs may have different graphs.
    \item Interventionally equivalent SCMs may have different graphs.
  \end{itemize}
\end{frame}

\begin{frame}{Example}
  \begin{Eg}
    \begin{itemize}
      \item Consider an SCM $M$
  with endogenous variables $X,Y$ and exogenous random variable $W$ and structural equations
    $$X = -X^3 + X + W, \qquad Y = \sin(X) - W$$
  \item Its graph is $G(M)$: \\
  \centerline{\begin{tikzpicture}
    \node[var] (W) at (0,0) {$W$};
    \node[var] (X) at (1.5,0) {$X$};
    \node[var] (Y) at (3,0) {$Y$};
    \draw[arr] (W) -- (X);
    \draw[arr,bend left] (W) edge (Y);
    \draw[arr] (X) -- (Y);
  \end{tikzpicture}}
\item Note: $X$ is not a parent of itself, because
  its structural equation is equivalent to $X = \sqrt[3]{W},$
  where $X$ does not appear on the r.h.s..
\item
  When considering exogenous random variables as latent, we can consider the marginalized graph $G(M)^{\setminus W}$:\\

  \medskip
  \centerline{\begin{tikzpicture}
    \node[var] (X) at (1.5,0) {$X$};
    \node[var] (Y) at (3,0) {$Y$};
    \draw[biarr,bend left] (X) edge (Y);
    \draw[arr] (X) -- (Y);
  \end{tikzpicture}}
    \end{itemize}
  \end{Eg}
\end{frame}

\begin{frame}{Example: Self-cycle}
\begin{Eg}
  The graph of an SCM with structural equation
  $$X = W X + c$$
  ($X$ endogenous, $W$ exogenous random) has a \emph{self-cycle} at $X$:\\
  \centerline{\begin{tikzpicture}
    \node[var] (W) at (0,0) {$W$};
    \node[var] (X) at (1.5,0) {$X$};
    \draw[arr] (W) -- (X);
    \draw[arr,in=30,out=-30,looseness=5] (X) edge (X);
  \end{tikzpicture}}
\end{Eg}
Self-cycles are a warning sign of complications with respect to solvability.
They are of no concern when restricting to the class of simple SCMs.
\begin{Prp}
  For $j \in V$, we have that there is a self-cycle $j \to j$ in $G(M)$ if and only if
  $M_{\doit(V\setminus \{j\})}$ is \emph{not} uniquely solvable.
  In particular, graphs of simple SCMs have no self-cycles.
\end{Prp}
\end{frame}

\begin{frame}{Compatibility of operations on SCMs and graphs}
  \begin{Prp}
  \begin{itemize}
    \item Hard interventions: for $T \subseteq J \cup V \cup W$,
      $$G(M_{\doit(T)}) = G(M)_{\doit(T)}.$$
    \item Marginalizations: If $M$ is simple, then for $L \subseteq V$,
      $$G(M_{\setminus L}) \subseteq G(M)^{\setminus L}.$$
  \end{itemize}
\end{Prp}
\begin{proof}
Try proving the first statement yourself!
The second statement is more involved. For a proof by induction, see the lecture notes.
\end{proof}
\end{frame}

\begin{frame}{Acyclic SCMs}
\begin{Def}
  An SCM $M$ is called \emph{acyclic} if its graph $G(M)$ is acyclic.
\end{Def}
  While most literature only considers acyclic SCMs (or CBNs), we propose to
  use simple SCMs instead if one cannot rule out a priori the possibility of causal cycles.
\begin{Prp}
  Acyclic SCMs are simple.
\end{Prp}
\begin{proof}
  Main idea: Choose a topological ordering and solve the system of structural equations in that order.
  For details, see lecture notes.
\end{proof}
\end{frame}
